% 张三慧《大学物理学》第4版V2第25章《几何光学》
% 源教材扫描PDF物理页341--343（印刷324--326）；配套辅导S PDF338--348。
% 图25.44--25.47及相应解答光路图需专家重绘复核。
\section*{第25章\quad 几何光学}
\subsection*{章末习题25.1--25.21}
\paragraph{习题 25.1}
高 \(8.0\,\mathrm m\) 的路灯下，
身高 \(1.70\,\mathrm m\) 的人以
\(1.5\,\mathrm{m/s}\) 的速度在水平道路上远离路灯。
求其头顶影子在地面上的移动速率。
\paragraph{习题 25.2}
泳者在岸边用手指和拇指围成的小孔观察
掉进泳池的眼镜。小孔离前胸的水平距离
\(30\,\mathrm{cm}\)、高于眼睛 \(30\,\mathrm{cm}\)，
离水面高 \(1.6\,\mathrm m\)，池深
\(2.0\,\mathrm m\)，水折射率 \(1.33\)。
求眼镜在水中距泳池竖直壁的水平距离。
\paragraph{习题 25.3}
空气中波长\emph{580 nm}的黄光，
以入射角45°射到金刚石表面，
折射角为17.0°。
求折射率、黄光在金刚石内的频率、
波长及速度。
\textit{按主教材580nm计算；配套S公式例值采用589nm需另列差异。}
\paragraph{习题 25.4}
如图25.44用作图法求折射线：
在入射线上取一点作圆心，
依折射率 \(n_1:n_2\) 作辅助圆，
借平行线和法线构造折射线。
从图中正弦定律证明
\(n_1\sin\theta_1=n_2\sin\theta_2\)。
\paragraph{习题 25.5}
红光 \(\lambda=633\,\mathrm{nm}\) 对某玻璃折射率
\(n_g=1.52\)。
求其从玻璃射向空气和从玻璃射向水
（\(n_w=1.33\)）时的全反射临界角。
\paragraph{习题 25.6}
空气中折射率1.52的玻璃表面垂直入射一束光。
不计吸收，求表面反射光强占入射光强百分之几。
\paragraph{习题 25.7}
对称玻璃三棱镜顶角 \(A\)，
光从一侧入射并对称地从另一侧出射，
其入射与出射方向的偏向角为 \(\delta\)
（图25.45）。
证明折射率
\(n=\sin((A+\delta)/2)/\sin(A/2)\)；
说明对称路线的偏向角为最小偏向角。
\paragraph{习题 25.8}
凹面镜球面曲率半径 \(R=30\,\mathrm{cm}\)。
物体分别置于镜前 \(s=25\,\mathrm{cm}\)
和 \(10\,\mathrm{cm}\)。
求像的位置、正倒、虚实、横向放大率，
并作成像光路图。
\paragraph{习题 25.9}
牙科小反射镜置于牙齿前 \(2.0\,\mathrm{cm}\)，
得到正立放大 \(5.0\) 倍的像。
求反射镜凹凸类型以及曲率半径。
\paragraph{习题 25.10}
图25.46：高40cm、宽20cm的暗箱底面
有曲率半径 \(40\,\mathrm{cm}\) 的凹镜，
箱顶部有与水平成45°的平面镜。
在凹镜竖直光轴上距其顶点30cm的位置，
水平拉住高3cm的小玉佛并照亮。
求最终像的位置、虚实、正倒、尺寸、
观察眼睛位置（考虑凹镜与平面镜两次反射）。
\paragraph{习题 25.11}
眼模型中，角膜为球面折射界面，
其前后材料的折射率分别为1.0和1.4，
其余折射只发生在该角膜。
角膜顶点到视网膜距离2.60cm。
(1) 平行光聚焦于视网膜时所需曲率半径；
(2) 将角膜前25.0cm的物体聚焦于视网膜
所需曲率半径。
\paragraph{习题 25.12}
球形鱼缸直径40cm，鱼位于水平半径的中点，
从缸外观察。
水折射率取1.33。求鱼的像相对鱼缸壁的位置、
虚实与正倒、横向放大倍数。
\paragraph{习题 25.13}
重火石玻璃在紫光 \(\lambda=410\,\mathrm{nm}\)
折射率1.698，在红光 \(\lambda=660\,\mathrm{nm}\)
折射率1.662。
两球面曲率半径大小均为20cm的双凸薄透镜，
求哪种颜色的焦点距镜更近，以及两焦点距离差。
\paragraph{习题 25.14}
图25.47：太阳光经球形水滴折射、一次内反射、
再折射后形成主虹。(1) 由光路证明
\(\alpha=4\theta_2-2\theta_1\)；
(2) 由极值条件 \(d\alpha/d\theta_1=0\)
证明 \(\cos^2\theta_{1c}=(n_w^2-1)/3\)；
(3) 红光、紫光折射率为1.333和1.342，
分别求极值入射角与出射折返角，解释虹内外颜色次序。
\paragraph{习题 25.15}
折射率 \(n=1.60\) 的平凸薄透镜，
球面半径24cm。
物体位于透镜前(1)120cm、(2)80cm、
(3)60cm、(4)20cm时，求像距、
正倒、虚实及大小，并作(3)(4)的光路图。
\paragraph{习题 25.16}
蜡烛与白屏相距 \(L\)；
在两者之间移动焦距 \(f\) 的凸透镜，
有两处能在屏上成清晰像，两位置相距 \(d\)。
证明 \(f=(L^2-d^2)/(4L)\)，
比较两次成像的放大率，并证明要有两位置
必满足 \(L>4f\)。
\paragraph{习题 25.17}
投影幻灯片用凸透镜焦距20.0cm，
投影屏距透镜5.00m。
求幻灯片应距透镜多远、如何放置以在屏上
得到正立放大像，以及屏上面积与幻灯面积之比。
\paragraph{习题 25.18}
两个共轴凸透镜焦距分别10cm与8cm，
沿水平光轴相距10cm。
在较近的第一透镜前15cm处放高1.5cm的小玉佛。
求经过两镜后最终像的位置、虚实、正倒及高度。
\paragraph{习题 25.19}
焦距 \(f_1,f_2\) 的两片薄透镜紧贴共轴，
证明等效焦距 \(f\) 满足 \(1/f=1/f_1+1/f_2\)。
\paragraph{习题 25.20}
美国Yerkes折射望远镜物镜直径1.02m、
物镜焦距19.5m、目镜焦距10cm；
Palomar反射望远镜物镜口径5.1m、
物镜焦距16.8m、目镜焦距1.25cm。
求两台望远镜的角放大率和波长550nm时的
最小分辨角。
\paragraph{习题 25.21}
一望远镜主物镜为口径10m、焦距20m的凹面镜
\(M_1\)，其前14m放球面镜 \(M_2\)。
欲使无限远星体的像形成于 \(M_1\) 后方4m，
求 \(M_2\) 的焦距及凹凸类型。
若目镜为焦距20cm的凸透镜，求其放置位置、
望远镜角放大率，以及星光波长590nm时的瑞利分辨角。
